\chapter{2021年「大学物理2」\titlelogo{https://sci.hdu.edu.cn}{HDU}期中试题}
\date{2021年11月20日}{大学物理教学团队}{A0715012}
{未央物理讲师 \href{https://qm.qq.com/q/UPbGudx8cK}{\textbf{Axia}}}
{\qrcode{https://qm.qq.com/q/asJiytD5pC}}{\href{https://qm.qq.com/q/asJiytD5pC}{HDU 物理营}}
{\qrcode{http://weixin.qq.com/r/mhKYgFTEMQlOrRDj90eI}}{\href{http://weixin.qq.com/r/mhKYgFTEMQlOrRDj90eI}{未央学社公众号}}

\section{选择题（每题3分，共27分）}
\begin{choice}{C}[单摆]
    把单摆摆球从平衡位置向位移正方向拉开，使摆线与竖直方向成一微小角度$\theta$，然后由静止放手任其振动，从放手时开始计时.若用余弦函数表示其运动方程，则该单摆振动的初相为
    \begin{tasks}(4)
        \task $\pi$
        \task $\frac{\pi}{2}$
        \task $0$
        \task $\theta$
    \end{tasks}
\end{choice}

\begin{choice}{B}[多普勒效应]
    一机车汽笛频率为$750\mathrm{Hz}$，机车以时速$90$公里远离静止的观察者，观察者听到声音的频率是（空气中声速$340\mathrm{m/s}$）
    \begin{tasks}(4)
        \task $810\mathrm{Hz}$
        \task $699\mathrm{Hz}$
        \task $805\mathrm{Hz}$
        \task $695\mathrm{Hz}$
    \end{tasks}
\end{choice}
\begin{solution}
    已知多普勒效应观察者（Observer）和发射源（Source）的的频率关系为
    \[\nu=\frac{u\pm v_o}{u\mp v_s}\nu_0\]
    \color{solutioncolor}
    $v_o$为观察者速度，接近为$+$，远离为$-$；$v_s$为发射源速度，接近为$-$，远离为$+$.观察者静止，其所听频率为
    \[\nu=\frac{340}{340+25}\times 750\mathrm{Hz}\approx 699\mathrm{Hz}\]
    
    \sokka{B}
\end{solution}

\begin{choice}{C}[光程和光程差]
    在相同的时间内，一束波长为$\lambda$的单色光在空气中和在玻璃中
    \begin{tasks}(2)
        \task 传播的路程相等，走过的光程相等
        \task 传播的路程相等，走过的光程不相等
        \task 传播的路程不相等，走过的光程相等
        \task 传播的路程不相等，走过的光程不相等
    \end{tasks}
\end{choice}
\begin{solution}*
    光程的定义：在相同时间内光线在真空中传播的距离.题目中光传播时间相同，故光程相等；又因为光在两种介质中的传播速度不同，所以在相同的时间内传播的路程不相等.\sokka{C}
\end{solution}

\begin{choice}{B}[双缝干涉]
    在双缝干涉实验中，为使屏上的干涉条纹间距变大，可采取的办法是
    \begin{tasks}(4)
        \task 使屏靠近双缝
        \task 使两缝的间距变小
        \task 把两缝的宽度调窄
        \task 改用波长短的单色光
    \end{tasks}
\end{choice}
\begin{solution}*
    已知双缝干涉条纹间距公式$\Delta x=\frac{\lambda L}{d}$.要使$\Delta x$变大，可增大波长$\lambda$、屏幕与双缝的距离$L$或减小双缝间距$d$.\sokka{B}
\end{solution}

\begin{paracol}{2}
    \begin{choice}{B}[双缝干涉]
        在双缝干涉实验中，屏幕$E$上的$P$点是明条纹.若将缝$S_2$盖住，并在$S_1S_2$连线的垂直平分面处放一高折射率反射面$M$，如图所示.则此时
    \begin{tasks}(2)
        \task $P$点仍为明条纹
        \task $P$点为暗条纹
        \task 不能确定$P$点是明纹还是暗纹
        \task 无干涉条纹
    \end{tasks}        
    \end{choice}
    \switchcolumn\centering
    \vfill
    \begin{tikzpicture}[decoration={markings,mark=between positions .2 and .8 step 18mm with {\arrow{stealth}}}]
        \draw [very thick] (0,1.5)--(0,-1.5);
        \draw [thick,postaction=decorate] (-1,0)--(0,0.5)--(2,0.8);
        \draw [thick] (0,0.5)--(0.8,0);
        \draw [thick,postaction=decorate] (-1,0)--(0,-0.5)--(2,0.8);
        \draw [very thick] (0,0)--(1.8,0);
        \draw [very thick] (2,1.8)--(2,-1.8);
        \fill [pattern=north west lines] (0,0) rectangle (1.8,-0.2);
        \fill [pattern=north east lines] (2,1.8) rectangle (2.2,-1.8);
        \node [anchor=east] at (-1,0) {$S$} node [anchor=south east] at (0,0.5) {$S_1$} node [anchor=north east] at (0,-0.5) {$S_2$};
        \node [anchor=north] at (0.8,-0.2) {$M$} node [anchor=west] at (2.2,0.8) {$P$} node [anchor=north] at (2.1,-1.8) {$E$};
        \node at (-1,0) {$\times$};
    \end{tikzpicture}
    \vfill
\end{paracol}
\begin{solution}*
    $S_1MP$、$S_2MP$长度相等，但平面镜使在反射中一条光路发生半波损失，两条光路的相位差变化$\pi$，所以$P$点由原来的明纹变为暗纹.\sokka{B}
\end{solution}

\begin{choice}{D}[增透膜]
    在照相机镜头的玻璃片上均匀镀有一层折射率$n$小于玻璃的介质薄膜，以增强某一波长$600\mathrm{nm}$的透射光能量.假设光线垂直入射，则介质膜的最小厚度应为
    \begin{tasks}(4)
        \task $\frac{600}{n}\mathrm{nm}$
        \task $\frac{300}{n}\mathrm{nm}$
        \task $\frac{200}{n}\mathrm{nm}$
        \task $\frac{150}{n}\mathrm{nm}$
    \end{tasks}
\end{choice}
\begin{solution}
    \begin{itemize}
        \item 需减少反射光的干涉作用，使反射光干涉相消.即$\delta=2ne=\frac{2k+1}{2}\lambda$，解得$e_{\min}=\frac{\lambda}{4n}=\frac{150}{n}\mathrm{nm}$.
        \item 也可使透射光干涉相长，即$\delta=2en+\frac{\lambda}{2}=k\lambda$，解得$e_{\min}=\frac{\lambda}{4n}=\frac{150}{n}\mathrm{nm}$.
    \end{itemize}
\end{solution}

\begin{choice}{D}[牛顿环]
    牛顿环干涉装置上平凸透镜在垂直于平板玻璃的方向上，逐渐向上平移（离开玻璃板）时，反射光形成的干涉条纹的变化情况是
    \begin{tasks}(2)
        \task 环纹向边缘扩散，环数不变
        \task 环纹向边缘扩散，环数增加
        \task 环纹向中心靠拢，环数增加
        \task 环纹向中心靠拢，环数不变
    \end{tasks}
\end{choice}
\begin{solution}*
    对于某条环，其光程差是确定的，所以环数不变；向中心靠拢光程差减小，可抵消透镜上移时导致的光程差增大.\sokka{D}
\end{solution}

\begin{choice}{A}[迈克尔逊干涉仪]
    在迈克尔逊干涉仪的一条光路中，放入一折射率为$n$，厚度为$d$的透明薄片，放入后，这条光路的光程改变了
    \begin{tasks}(4)
        \task $2\ab(n-1)d$
        \task $\ab(n-1)d$
        \task $2\ab(n-1)d+\frac{\lambda}{2}$
        \task $nd$
    \end{tasks}
\end{choice}
\begin{solution}*
    迈克尔逊干涉仪光路上的某点光碰到反射镜后会再次经过此处，故原来的光程为$2d$；现在加入了透明薄片，使得这里的光程为$2nd$，故光程差为$2\ab(n-1)d$.\sokka{A}
\end{solution}

\begin{choice}{C}[弗琅禾费衍射]
    波长为$\lambda$的单色平行光垂直入射到一狭缝上，若第一级暗纹的位置对应的衍射角为$\theta=\pm\frac{\pi}{6}$，则缝宽的大小为
    \begin{tasks}(4)
        \task $\frac{\lambda}{2}$
        \task $\lambda$
        \task $2\lambda$
        \task $3\lambda$
    \end{tasks}
\end{choice}
\begin{solution}*
    一级暗纹对应的方程为$d\sin\frac{\pi}{6}=1\cdot\lambda$，得$d=2\lambda$.\sokka{C}
\end{solution}

\section{填空题（共25分）}

\begin{problem}[简谐振动][3]
    用$40\mathrm{N}$的力拉一轻弹簧，可使其伸长$20\mathrm{cm}$.此弹簧下应挂\ans{$2$}$\mathrm{kg}$的物体，才能使其做简谐振动的周期为$T=0.2\pi$.
\end{problem}
\begin{solution}*
    弹簧劲度系数$k=\frac{40\mathrm{N}}{0.2\mathrm{m}}=200\mathrm{N/m}$；由振子振动周期$T=2\pi\sqrt{\frac{m}{k}}$得$m=\frac{kT^2}{4\pi^2}=2\mathrm{kg}$.
\end{solution}

\begin{paracol}{2}
\begin{problem}[简谐振动][4]
    两个简谐振动曲线如图所示，二者频率之比为$\nu_1:\nu_2=$\ans{$1:2$}，加速度最大值之比为$a_{1m}:a_{2m}$\ans{$1:4$}，初始速率之比为$v_{10}:v_{20}=$\ans{$1:2$}.
\end{problem}    
    \switchcolumn\centering
    \vfill
    \begin{tikzpicture}[scale=0.83]
        \draw [->,thick,](-0.5,0)--(3.5,0);
        \draw [->,thick](0,-1.2)--(0,1.5);
        \draw[thick,H1] (0,0) sin (0.5,1) cos (1,0) sin (1.5,-1) cos (2,0) sin (2.5,1) cos (3,0) ;
        \draw[thick,2号色] (0,0) sin (0.25,1) cos (0.5,0) sin (0.75,-1) cos (1,0) sin (1.25,1) cos (1.5,0) sin (1.75,-1) cos (2.0,0) sin (2.25,1) cos (2.5,0) sin (2.75,-1);
    \end{tikzpicture}
    \vfill
\end{paracol}
\begin{solution}
    \begin{itemize}
        \item 由图易知，$x_1$与$x_2$的周期之比为$2:1$，故两者频率之比为周期之比的倒数，即$\nu_1:\nu_2=1:2$；
        \item 加速度最大值$a_m=\omega^2 A\propto\omega^2\propto\nu^2$，所以$a_{1m}:a_{2m}=\omega_1^2:\omega_2^2=\nu_1^2:\nu_2^2=1:4$；
        \item $t=0$时质点恰好在平衡点，此时初始速率为最大速率，$v_{1m}:v_{2m}=\omega_1A:\omega_2A=\nu_1:\nu_2=1:2$.
    \end{itemize}
\end{solution}

\begin{paracol}{2}
    \begin{problem}[平面简谐波的波函数][4]
        设某时刻一横波波形曲线如图所示.
        \begin{enumerate}
            \item 试分别用矢量符号表示图中A、B、C、D、E、F、G、H、I质点在该时刻的运动方向.
            \item 画出四分之一周期后的波形曲线.
        \end{enumerate}
    \end{problem}
    \switchcolumn\centering
    \vfill
    \begin{tikzpicture}[scale=1.2,decoration={markings,
        mark=between positions 0 and 1 step 7.23mm
        with {\shade[ball color=darkgray](0,0) circle (0.05);}}]
        \draw [->,thick] (-0.5,0)--(3,0);
        \draw [->,thick] (0,-1.35)--(0,1.35);
        \draw[thick,H1,postaction=decorate] (0,0) sin (0.6,1) cos (1.2,0) sin (1.8,-1) cos (2.4,0);
        \node [anchor=south east] at (0,0) {$A$} node [anchor=south east] at (0.2,0.5) {$B$} node [anchor=south] at (0.6,1.0) {$C$} node [anchor=east] at (1.0,0.5) {$D$} node [anchor=north east] at (1.2,0) {$E$} node [anchor=north east] at (1.4,-0.5) {$F$} node [anchor=north] at (1.8,-1.0) {$G$} node [anchor=east] at (2.2,-0.5) {$H$} node [anchor=south] at (2.4,0) {$I$};
        \draw [very thick,2号色,->] (0,0)--(0,-0.5);
        \draw [very thick,2号色,->] (0.23,0.56)--(0.23,0.06);
        \draw [very thick,2号色,->] (0.97,0.56)--(0.97,1.06);
        \draw [very thick,2号色,->] (1.2,0)--(1.2,0.5);
        \draw [very thick,2号色,->] (1.43,-0.56)--(1.43,-0.06);
        \draw [very thick,2号色,->] (2.17,-0.56)--(2.17,-1.06);
        \draw [very thick,2号色,->] (2.4,0)--(2.4,-0.5);
        \draw [very thick,dash dot,H2] (0.0,-1) cos (0.6,0) sin (1.2,1) cos (1.8,0) sin (2.4,-1);
    \end{tikzpicture}
    \vfill
\end{paracol}

\begin{problem}[弹簧振子][3]
    一作简谐振动的振动系统，振子质量为$2\mathrm{kg}$，系统振动频率为$1000\mathrm{Hz}$，振幅为$0.5\mathrm{cm}$，则其振动能量为$\ans{987\mathrm{J}}$
\end{problem}
\begin{solution}*
    由$\omega=2\pi\nu=\sqrt{\frac{k}{m}}$得$k=4\pi^2m\nu^2$.代入振动能量表达式得$E=\frac{1}{2}kA^2=2\pi^2m\nu^2A^2\approx 987\mathrm{J}$.
\end{solution}

\begin{problem}[平面简谐波的波函数][5]
    一平面简谐波沿$Ox$轴\textbf{\textcolor{H1}{负}}方向传播，波长为$\lambda$，若位于$x=-L$的$P$处质点的振动方程为$y_p=A\cos\ab(2\pi\nu t+\frac{\pi}{2})$，则该波的表达式为\ans{$y=A\cos\ab[2\pi\ab(\nu t+\frac{x+L}{\lambda})+\frac{\pi}{2}]$}；$P$处质点\ans{$t_1+\frac{L}{\lambda\nu}+\frac{k}{\nu}$}时刻的振动状态与$O$处质点$t_1$时刻的振动状态相同.
\end{problem}
\begin{solution}
    设波函数为$y=A\cos\ab[2\pi\nu\ab(t+\frac{x}{\lambda\nu})+\varphi_0]$.将$P$点坐标$x=-L$代入，得
    \[y_p=A\cos\ab[2\pi\nu\ab(t-\frac{L}{\lambda\nu})+\varphi_0]\]
    \color{solutioncolor}
    将上式与$P$点振动方程比较，得初相$\varphi_0=\frac{\pi}{2}+\frac{2\pi L}{\lambda}$.将$\varphi_0$代入波函数中，得该波的表达式
    \[y=A\cos\ab[2\pi\nu t+\frac{2\pi\ab(x+L)}{\lambda}+\frac{\pi}{2}]\]

    令$x=0$，得到$t_1$时刻$O$点的振动方程$y_o(t_1)=A\cos\ab(2\pi\nu t_1+\frac{2\pi L}{\lambda}+\frac{\pi}{2})$. 由$P$处质点在$t$时刻的振动状态与$O$处质点在$t_1$时刻的振动状态相同得
    \[y_p(t)=A\cos\ab(2\pi\nu t+\frac{\pi}{2})=y_o(t_1)A\cos\ab(2\pi\nu t_1+\frac{2\pi L}{\lambda}+\frac{\pi}{2}+2k\pi)\]

    解得$t=t_1+\frac{L}{\lambda\nu}+\frac{k}{\nu}$.
\end{solution}

\begin{problem}[迈克尔逊干涉仪][3]
    用迈克尔逊干涉仪测微小的位移，若入射光波波长$\lambda=628.9\mathrm{nm}$，当动臂反射镜移动时，干涉条纹移动了$2048$条，反射镜移动的距离$d=\ans{0.644\mathrm{mm}}$.
\end{problem}
\begin{solution}*
    移动带来的光程差满足$\delta=2d=N\lambda$，由此得$d=N\lambda/2=0.644\mathrm{mm}$.
\end{solution}

\begin{problem}[弗琅禾费衍射][3]
    平行单色光垂直入射在缝宽为$a=0.15\mathrm{mm}$的单缝上，缝后有焦距为$f=400\mathrm{mm}$的凸透镜，在其焦平面上放置观察屏幕.现测得屏幕上中央明条纹两侧的两个第三级暗纹之间的距离为$8\mathrm{mm}$，则入射光的波长为$\lambda=\ans{500\mathrm{nm}}$.
\end{problem}
\begin{solution}
    设$\pm 3$级条纹与光轴的夹角为$\theta$，已知$\pm 3$级条纹的坐标为$\pm 4\mathrm{mm}$，由几何关系得
    \[a\sin\theta=3\lambda,\ \frac{3\lambda}{a}=\sin{\theta}\approx\tan\theta=\frac{4\mathrm{mm}}{f}=0.01\]
    得入射光波长为$\lambda=a\theta/3=500\mathrm{nm}$.
\end{solution}

\section{计算题（共48分）}

\begin{problem}[简谐振动][10]
    一轻弹簧下悬挂$m_0=100\mathrm{g}$砝码时，弹簧伸长$8\mathrm{cm}$.现在这根弹簧下端悬挂$m=250\mathrm{g}$的物体构成弹簧振子.将物体从平衡位置向下拉动$4\mathrm{cm}$，并给以向上$21\mathrm{cm/s}$的初速度（令这时$t=0$）.选$x$轴向下，求振动方程的数值式.
\end{problem}
\begin{solution}
    \begin{itemize}
        \item 弹簧的劲度系数、角频率分别为$k=m_0g/\Delta x_0=12.5\mathrm{N/m},\ \omega=\sqrt{k/m}=7.07\mathrm{s}$.\point{3}
        \item 由旋转矢量法得系统振幅、初相为$A=\sqrt{x_0^2+\ab(\frac{v_0}{\omega})^2}=4.98\mathrm{cm},\ \varphi_0=\arccos{\frac{x_0}{A}}=0.20\pi$.\point{4}
        \item 振动方程为$y=4.98\cos\ab(7.07t+0.20\pi)$.\point{3}
    \end{itemize}
    \end{solution}

\begin{problem}[简谐振动的合成][5]
    两个同方向简谐振动的振动方程分别为
    \[x_1=5\times 10^{-2}\cos\ab(10t+\frac{3}{4}\pi)\ (\mathrm{SI}),\ x_2=6\times 10^{-2}\cos\ab(10t+\frac{1}{4}\pi)\ (\mathrm{SI})\]

    求合振动方程.
\end{problem}
\begin{solution}
    \begin{itemize}
        \item 合振动振幅为$A=\sqrt{A_1^2+A_2^2+2A_1A_2\cos\ab(\varphi_1-\varphi_2)}=7.8\times 10^{-2}$.\point{2}
        \item 合振动相位$\varphi=\arctan{\frac{A_1\sin\varphi_1+A_2\sin\varphi_2}{A_1\cos{\varphi_1}+A_2\cos{\varphi_2}}}=0.47\pi$.\point{2}
        \item 两者振动频率相同，所以合振动频率与其相同.故合振动方程为$x=7.8\times 10^{-2}\cos\ab(10t+0.47\pi)$.\point{1}
    \end{itemize}
\end{solution}

\begin{problem}[平面简谐波的波函数][5]
    一振幅为$10\mathrm{cm}$，波长为$200\mathrm{cm}$的简谐横波，沿着一条很长的水平的绷紧弦从左向右行进，波速为$100\mathrm{cm/s}$.取弦上一点为坐标原点，$x$轴指向右方，在$t=0$时原点处质点从平衡位置开始向位移负方向运动.求以SI单位表示的波动表达式（用余弦函数）及弦上任一点的最大振动速度.
\end{problem}
\begin{solution}
    由题意得，角频率$\omega=2\pi\frac{u}{\lambda}=\pi\mathrm{rad/s}$；$t=0$时$x=0$，$v<0$，所以初相$\varphi=\frac{\pi}{2}$.故波动表达式为\point{1}
    \[y=0.1\cos\ab[\pi\ab(t-x)+\frac{\pi}{2}]\ (\mathrm{SI})\eqno\point{2}\]
    最大振动速度$v_m=\omega A=0.1\pi\mathrm{m/s}$.\point{1}
\end{solution}

\begin{paracol}{2}
    \begin{problem}[双缝干涉][10]
        双缝干涉实验装置如图所示, 双缝与屏之间的距离$D=150\mathrm{cm}$, 两缝之间的距离$d=0.50\mathrm{mm}$, 用波长$\lambda=600\mathrm{nm}$的单色光垂直照射双缝.
        \begin{enumerate}
            \item 求原点$O$（零级明条纹所在处）上方第五级明条纹的坐标
            \item 如果用厚度$l=1.0\times 10^{-2}\mathrm{mm}$，折射率$n=1.58$的透明薄膜覆盖在图中的$S_1$缝后面，求上述第五级明条纹的坐标$x^{\prime}$.
        \end{enumerate}
    \end{problem}
    \switchcolumn\centering
    \vfill
    \begin{tikzpicture}[scale=0.48]
        \draw [dash dot] (-0.2,0)--(6.2,0);
        \draw (0,1.5)--(0,3);
        \draw (0,-1)--(0,1);
        \draw (0,-1.5)--(0,-3);
        \draw [2号色,thick] (0,1.25)--(6,2.25);
        \draw [2号色,thick] (0,-1.25)--(6,2.25);
        \draw (6,3)--(6,-3);
        \draw [<->] (0,-2.5)--(6,-2.5) node [anchor=south,midway] {$D$};
        \node [anchor=south west] at (0,1.25) {$S_1$} node [anchor=north west] at (0,-1.25) {$S_2$};
        \node [anchor=east] at (0,0) {$d$};
        \draw[->] (-1.25,1.25)--(-0.2,1.25);
        \draw[->] (-1.25,-1.25)--(-0.2,-1.25);
    \end{tikzpicture}
    \vfill
\end{paracol}
\begin{solution}
    \begin{enumerate}
        \item $x_5=\frac{5\lambda D}{d}=9\mathrm{mm}$.\point{4}
        \item 根据明纹条件，此时的光程差为
        \[\delta=r_2-r_1-\ab(n-1)l=\frac{x_5^{\prime}d}{D}-\ab(n-1)l=k\lambda\eqno\point{4}\]
        得$x_5^{\prime}=\ab[k\lambda+\ab(n-1)l]\frac{D}{d}=26.4\mathrm{mm}$.\point{2}
    \end{enumerate}
\end{solution}

\begin{problem}[简谐振动的合成][5]
    三个频率相同、振动方向相同（垂直纸面）的简谐波，在传播过程中在$O$点相遇；若三个简谐波各自单独在$S_1$、$S_2$和$S_3$的振动方程分别为$y_1=A\cos\ab(\omega t+\frac{1}{2}\pi)$，$y_2=A\cos\omega t$和$y_3=2A\cos\ab(\omega t-\frac{1}{2}\pi)$；且$\overline{S_2O}=4\lambda$，$\overline{S_1O}=\overline{S_3O}=5\lambda$，求$O$点的合振动方程（设传播过程中各波振幅不变）.
\end{problem}
\begin{solution}
    三个波源到$O$点的距离都是波长的整数倍，无需考虑传播到$O$点过程带来的相位差.故$O$点的合振动为\point{1}
    \[\widetilde{y}=\widetilde{y_1}+\widetilde{y_2}+\widetilde{y_3}=A\e^{\i\cdot 0}+A\e^{\i\frac{\pi}{2}}+2A\e^{-\i\frac{\pi}{2}}=\sqrt{2}A\e^{-\i\frac{\pi}{4}}\to y=\sqrt{2}A\cos\ab(\omega t-\frac{\pi}{4})\eqno\point{4}\]
\end{solution}

\begin{problem}[牛顿环][8]
    牛顿环装置透镜凸表面的曲率半径是$R=400\mathrm{cm}$.用某单色平行光垂直入射，观察反射光形成的牛顿环，测得第5个明环的半径是$0.30\mathrm{cm}$.
    \begin{enumerate}
        \item 求入射光的波长.
        \item 求以透镜中心为圆心在半径为$1\mathrm{cm}$的范围内可观察到的明环数目.
    \end{enumerate}
\end{problem}
\begin{solution}
    \begin{enumerate}
    \item 由明环半径公式$r_k=\sqrt{\frac{\ab(2k-1)\lambda R}{2}}$得入射光的波长$\lambda=\frac{2r^2}{\ab(2k-1)R}=500\mathrm{nm}$.\point{5}
    \item 由明环半径公式，代入$r=1\mathrm{cm}$得$k=\left\lfloor\frac{r^2}{\lambda R}+\frac{1}{2}\right\rfloor=50$.\point{3}
    \end{enumerate}
\end{solution}

\begin{problem}[弗琅禾费衍射][5]
在单缝的弗琅禾费衍射中，缝宽$a=0.100\mathrm{mm}$，平行光垂直入射在单缝上，波长$\lambda=500\mathrm{nm}$，会聚透镜的焦距$f=1.00\mathrm{m}$.求中央亮纹旁的第一个亮纹的宽度$\Delta x$.
\end{problem}
\begin{solution}
    求出其两侧暗纹位置后作差即可求出亮纹宽度.
\begin{itemize}
    \item 一级暗纹
    $a\sin\theta_1=\lambda,\ x_1=f\tan\theta_1\approx f\sin\theta_1=\frac{\lambda f}{a}$。\point{2}
    \item 二级暗纹
    $a\sin\theta_2=2\lambda,\ x_2=f\tan\theta_2\approx f\sin\theta_2=\frac{2\lambda f}{a}$.\point{2}
    \item 亮纹宽度
    $\Delta x=x_2-x_1=\frac{\lambda f}{a}=5\mathrm{mm}$.\point{1}
\end{itemize}
\end{solution}